
\let\negmedspace\undefined
\let\negthickspace\undefined
\documentclass[journal]{IEEEtran}
\usepackage[a5paper, margin=10mm, onecolumn]{geometry}
%\usepackage{lmodern} % Ensure lmodern is loaded for pdflatex
\usepackage{tfrupee} % Include tfrupee package

\setlength{\headheight}{1cm} % Set the height of the header box
\setlength{\headsep}{0mm}     % Set the distance between the header box and the top of the text

\usepackage{gvv-book}
\usepackage{gvv}
\usepackage{cite}
\usepackage{amsmath,amssymb,amsfonts,amsthm}
\usepackage{algorithmic}
\usepackage{graphicx}
\usepackage{textcomp}
\usepackage{xcolor}
\usepackage{txfonts}
\usepackage{listings}
\usepackage{enumitem}
\usepackage{mathtools}
\usepackage{gensymb}
\usepackage{comment}
\usepackage[breaklinks=true]{hyperref}
\usepackage{tkz-euclide} 
\usepackage{listings}
% \usepackage{gvv}                                        
\def\inputGnumericTable{}                                 
\usepackage[latin1]{inputenc}                                
\usepackage{color}                                            
\usepackage{array}                                            
\usepackage{longtable}                                       
\usepackage{calc}                                             
\usepackage{multirow}                                         
\usepackage{hhline}                                           
\usepackage{ifthen}                                           
\usepackage{lscape}
\begin{document}

\bibliographystyle{IEEEtran}
\vspace{3cm}

\title{12.775}
\author{EE25BTECH11033 - Kavin}
% \maketitle
% \newpage
% \bigskip
{\let\newpage\relax\maketitle}

\renewcommand{\thefigure}{\theenumi}
\renewcommand{\thetable}{\theenumi}
\setlength{\intextsep}{10pt} % Space between text and floats
\textbf{Question}:\\
Choose solution set $S$ corresponding to the systems of two equations
    \begin{align*}
    x - 2y + z &= 0 \\
    x - z &= 0
    \end{align*}
    \begin{multicols}{2}
    \begin{enumerate}
        \item $S = \left\{ \alpha \myvec{ 1 \\ 1 \\ 1 } \middle| \alpha \in \mathbb{R} \right\}$
        \item $S = \left\{ \alpha \myvec{ 1 \\ 1 \\ 1 } + \beta \myvec{ 1 \\ 0 \\ 1 } \middle| \alpha, \beta \in \mathbb{R} \right\}$
        \item $S = \left\{ \alpha \myvec{ 1 \\ 1 \\ 1 } + \beta \myvec{ 2 \\ 1 \\ 2 } \middle| \alpha, \beta \in \mathbb{R} \right\}$
        \item $S = \left\{ \alpha \myvec{ 1 \\ 0 \\ 1 } \middle| \alpha \in \mathbb{R} \right\}$
    \end{enumerate}
\end{multicols}\\    
\bigskip

\textbf{Solution}:\\


We are given a homogeneous system of two linear equations with three variables ($x$, $y$, $z$). We need to find all vectors $\begin{myvec}{ x \\ y \\ z }\end{myvec}$ that satisfy both equations simultaneously.

We can solve this system using substitution.

\begin{itemize}
    \item \textbf{From the second equation:}
    \begin{align}
    x - z = 0 \implies x = z
    \end{align}
    
    \item \textbf{Substitute $x = z$ into the first equation:}
    \begin{align}
        (z) - 2y + z &= 0 \\
        2z - 2y &= 0 \\
        2z &= 2y \\
        z &= y
    \end{align}
    
    \item \textbf{Combine the results:}
    We have found that $x = z$ and $y = z$. Therefore, all three variables must be equal:
    \begin{align}
    x = y = z
    \end{align}
\end{itemize}


Since we have three variables ($x, y, z$) but only two independent equations, we have one free variable. We can express the solution set in terms of a single parameter, let's call it $\alpha \in \mathbb{R}$.

\begin{itemize}
    \item Let the free variable be $z = \alpha$.
    \item Since $y = z$, then $y = \alpha$.
    \item Since $x = z$, then $x = \alpha$.
\end{itemize}

The solution vectors are of the form:
\begin{align}
\begin{myvec}{ x \\ y \\ z }\end{myvec} = 
\begin{myvec}{ \alpha \\ \alpha \\ \alpha }\end{myvec}
\end{align}

We can factor out the parameter $\alpha$:
\begin{align}
\alpha \begin{myvec}{ 1 \\ 1 \\ 1 }\end{myvec}
\end{align}
The solution set $S$ is the set of all possible vectors of this form for any real number $\alpha$.

Therefore, the solution set is:
\begin{align}
S = \left\{ \alpha \begin{myvec}{ 1 \\ 1 \\ 1 }\end{myvec} \mid \alpha \in \mathbb{R} \right\}
\end{align}


\begin{figure}[H]
\begin{center}
\includegraphics[width=0.6\columnwidth]{figs/fig.png}
\end{center}
\label{fig:Fig1}
\end{figure} 
\end{document}



